\subsection{Global Spectral Capacity and the Continuum Limit}
  \label{subsec:spectral-capacity}

  The structural analysis developed in the previous sections identifies a
  restricted set of finite relational graphs capable of supporting admissible
  neutral generating sets.
  To compare such structures quantitatively, it is useful to introduce a
  global measure of admissibility.

  Let $L$ denote the relational Laplacian associated with the Cayley graph of
  a generating set $S$.
  If $\{\lambda_\rho\}$ denotes the set of Laplacian eigenvalues associated
  with irreducible sectors $\rho$, the admissibility envelope introduced in
  Section~\ref{subsec:gram-rigidity} implies that the maximal effective
  amplitude in sector $\rho$ is bounded by
  \[
    A_\rho^{\max} = \frac{c_\chi}{\sqrt{\lambda_\rho}} .
  \]

  This relation naturally leads to a notion of global spectral capacity.
  A simple aggregate measure can be defined by summing the admissible
  amplitudes over all sectors,
  \begin{equation}
    C = \sum_\rho A_\rho^{\max}.
  \end{equation}

  However, constructive alignment between sectors can artificially inflate
  this quantity.
  To account for this effect one may introduce a penalised capacity
  \[
    C_\alpha,
  \]
  in which correlated contributions are weighted by a parameter
  $\alpha \in [0,1]$.

  This framework allows different relational structures to be compared in
  terms of their ability to support admissible spectral modes.
  Binary polyhedral groups, in particular the icosahedral group $2I$, emerge
  as near-maximal configurations among finite groups with comparable
  valency.

  A complementary diagnostic is provided by the admissibility ratio
  \[
    R_p ,
  \]
  which measures the relative concentration of admissible spectral weight
  across the hierarchy of sectors.
  Along sequences of increasingly refined relational graphs,
  this ratio approaches a stable value characteristic of the underlying
  spectral geometry.

  These considerations suggest that the structures identified in the
  preceding sections should not be viewed as isolated algebraic cases.
  Rather, they form part of a broader sequence of relational graphs whose
  spectral properties progressively approximate those of the continuous
  group $SU(2)$.

  One may therefore consider chains of relational structures of the form
  \[
    Q_8 \rightarrow 2I \rightarrow \mathrm{LPS} \rightarrow SU(2),
  \]
  in which increasingly refined graphs approximate the spectral geometry of
  the continuous spin group.

  In this limit the discrete spectrum becomes dense and the admissibility
  filter progressively weakens.
  The bounded-flux constraint still imposes a local restriction on the
  amplitude of modes, but the spectral hierarchy itself approaches a
  continuous distribution.

  The discrete structures analysed in this chapter can thus be interpreted as
  finite approximations of a deeper spinorial framework.
  Their role is to reveal the structural mechanisms through which admissible
  spectral hierarchies emerge from bounded relational dynamics.
